\section{Strict categories, directed chains, and simplicial gluing}
\label{sec:foundations}

The pushout argument needs three foundations: directed chain maps must
give strict right cylinders, these cylinders must descend to finite
dimensions, and their nerve comparisons must survive gluing and filtered
colimits. We fix the conventions and prove the needed local statements
before constructing the relative prism.

We work with small objects in a fixed universe.  An $\omega$-category in
this paper is a strict globular $\omega$-category: its cells have globular
sources and targets, identities, and compositions along matching
lower-dimensional boundaries, satisfying the associativity, unit, and
interchange laws strictly.  An $n$-category, for finite $n$, has only
identity cells above dimension $n$.  The inclusion
$\nCat\hookrightarrow\omCat$ has a left adjoint $\trunc$.
Explicitly, $\trunc X$ retains the cells of $X$ in dimensions below $n$,
and its $n$-cells are the $n$-cells of $X$ modulo the congruence generated
by
\[
 s_n(z)=t_n(z)\qquad(z\text{ an }(n+1)\text{-cell of }X).
\]
Here the congruence is closed under every defined composition.  Its
generating pairs are parallel $n$-cells, so their identifications leave
every lower matching boundary unchanged; compositions therefore descend
to the quotient.  Cells above $n$ become identities.  Every strict functor from $X$ to an
$n$-category annihilates these relations, which proves the stated
reflection property.  We set $\tau_\infty=\id$ and use $\nCat=\omCat$
when $n=\infty$.

Write $\Disk{i}$ for the free strict $\omega$-category on one $i$-cell,
with all its globular faces.  Thus
$\Hom_{\omCat}(\Disk{i},X)$ is the set of $i$-cells of $X$,
$\Disk{0}$ is terminal, and $\Disk{1}$ is the category $0\to1$.
Let $\Orient{m}$ denote the $m$th oriental.  Its standard cosimplicial
structure defines
\[
 (\Street X)_m=\Hom_{\omCat}(\Orient{m},X)
              =\Hom_{\nCat}(\trunc\Orient{m},X).
\]
For $\omega$-categories we also write $N_\infty$.  The left adjoints are
given by the simplex colimits
\begin{equation}\label{eq:categorification-colimit}
 c_\infty K=\colim_{(\Delta[m]\to K)\in\Delta/K}\Orient{m},
 \qquad
 c_nK=\trunc c_\infty K
     =\colim_{(\Delta[m]\to K)\in\Delta/K}\trunc\Orient{m}.
\end{equation}
These formulas follow from the universal property of a colimit and the
degreewise definition of the nerve.  In particular, for a poset $E$,
\[
 \Uinf{E}=c_\infty\OrdN E,\qquad
 \U{E}=c_n\OrdN E=\trunc\Uinf{E}.
\]
The fixed dimension is suppressed in $U$: for a monotone map $f$,
write $U(f)=c_n\OrdN(f)$ and $U^\infty(f)=c_\infty\OrdN(f)$.
We use $\OrdN$ for the ordinary nerve.  On ordinary categories it agrees
with the Street nerve: a map from an oriental is determined by its
successive edges, because its triangular cells impose the usual
composition equations.  In particular $N_\infty\Disk{1}=\Delta[1]$.

\subsection{The directed tensor convention}

An augmented directed complex is a nonnegative chain complex $K$ of
abelian groups, with augmentation $\epsilon:K_0\to\mathbb Z$ and a
specified additive submonoid $K_p^+\subset K_p$ in each degree.
Its morphisms are chain maps preserving the augmentation and these
submonoids.  All complexes used below have a specified free basis, with
positive submonoids generated by the basis.  For such a complex write
$dc=d^+c-d^-c$ with disjoint positive supports.  A basis is
\emph{unital} when
$\epsilon((d^-)^p b)=\epsilon((d^+)^p b)=1$ for each basis element
$b$ of degree $p$, and is \emph{strongly loop-free} when the relations
\[
 a<b\quad\text{if }a\text{ occurs in }d^-b,
 \qquad
 b<a\quad\text{if }a\text{ occurs in }d^+b
\]
are contained in a strict partial order on the whole basis.
A complex with such a basis is called a \emph{strong Steiner complex}.

Steiner's functor $\nu$ sends augmented directed complexes to strict
$\omega$-categories.  Its cells retain the lower-dimensional boundary
data of a chain.  More concretely, a $p$-cell of $\nu K$ is a table
\begin{equation}\label{eq:steiner-cell-table}
 (x_0^-,x_0^+\mid x_1^-,x_1^+\mid\cdots\mid x_p^-,x_p^+),
 \qquad x_r^\pm\in K_r^+,
\end{equation}
such that
\[
 \epsilon x_0^- = \epsilon x_0^+ = 1,\qquad
 d x_r^- = d x_r^+ = x_{r-1}^+-x_{r-1}^-
 \quad(1\le r\le p),\qquad x_p^-=x_p^+.
\]
An identity appends a zero pair to the table.  A positive,
augmentation-preserving chain map acts on every entry, giving the
corresponding strict functor
\cite[Definitions~2.6 and~2.8, Proposition~2.7]{Steiner2004}.
For a basis element $b$ of degree $p$, its \emph{atom} has top pair
$(b,b)$ and entries
$\langle b\rangle_r^\pm=(d^\pm)^{p-r}b$ for $0\le r\le p$:
at each step, take the negative or positive part of the boundary.
For a strong Steiner complex these tables are cells, and $\nu K$ is
freely generated as a polygraph by the atoms of its basis.  In
particular, maps out of it are determined by the images of these atoms
\cite[\S3 before Definition~3.2, Theorem~6.1]{Steiner2004}.
The concrete complexes needed here are as follows:
\begin{enumerate}[label=(\roman*)]
\item
The normalized simplex complex $M[m]$ has degree-$p$ basis
$[i_0,\ldots,i_p]$ for $0\le i_0<\cdots<i_p\le m$, with
\[
 d[i_0,\ldots,i_p]
   =\sum_{j=0}^{p}(-1)^j[i_0,\ldots,\widehat{i_j},\ldots,i_p]
 \quad(p>0),\qquad \epsilon[i]=1.
\]
It is strong Steiner and $\nu M[m]=\Orient{m}$.  A monotone map of
ordinals acts by applying it to the tuple; a tuple with a repeated
entry is zero.  The atom of a nonempty tuple has its first and last
vertices as its $0$-source and $0$-target.
\item
The globe complex $G[i]$ has one top basis element in degree $i$ and
two basis elements, its source and target faces, in each degree below
$i$.  Its positive-degree differentials are target minus source, and
its vertices have augmentation one.  It is strong Steiner and
$\nu G[i]=\Disk{i}$.  For $i=0$ there is just one vertex.
\end{enumerate}
The strong Steiner properties and tensor closure follow from
\cite[Examples~3.8--3.10]{Steiner2004}. The disk identification follows
from \cite[Example~4.7 and Theorem~5.11]{Steiner2004}. For the oriental
identification and its cosimplicial action, including repeated entries,
see \cite[Theorem~3.2 and \S4 before Theorem~4.1]{Steiner2007}.

The tensor of augmented directed complexes uses the ordinary
nonnegative chain-complex tensor, with separate augmentation:
\begin{equation}\label{eq:tensor-differential}
 d(x\otimes y)=dx\otimes y+(-1)^{|x|}x\otimes dy,
 \qquad
 \epsilon(x\otimes y)=\epsilon(x)\epsilon(y)
 \quad(|x|=|y|=0).
\end{equation}
Its positive submonoid is generated by the tensors of positive
elements.  Tensor products of the indicated bases are the relevant
positive bases.

We use the biclosed Gray tensor product on \emph{all} strict
$\omega$-categories supplied by
\cite[Theorem~A.15 and \S A.16]{AraMaltsiniotis2020}.
It has unit $\Disk{0}$, preserves small colimits separately in each
variable, and has canonical natural identifications
\begin{equation}\label{eq:steiner-tensor}
 \nu K\otimes\nu L\ \cong\ \nu(K\otimes L)
 \qquad(K,L\text{ strong Steiner}).
\end{equation}
Naturality here is for every morphism between strong Steiner
complexes, including maps sending basis elements to sums or zero.
This is the completed extension theorem; the general extension is
not being inferred from the incomplete proof of
\cite[Theorem~7.3]{Steiner2004}.
In particular, preservation of colimits in this paper is a property
of the Gray tensor, not a claimed preservation property of $\nu$.

It follows that $\Disk{i}\otimes\Disk{1}$ is generated by the atoms
indexed by pairs of globe basis elements, of dimension the sum of
their dimensions.  The same holds for
$\Orient{m}\otimes\Disk{1}$.  Globe face inclusions and simplex face
inclusions identify the corresponding based subcomplexes, and hence
the corresponding atoms in these tensor products.  One way to see the
last assertion directly is to compute the positive and negative
boundaries of a basis element within the based subcomplex: they are
unchanged by its inclusion.  The atom construction therefore gives
the same cell in the larger complex.

Throughout, the interval is the \emph{right} tensor factor.  Write
$e_0,e_1$ for the degree-zero basis elements of $M[1]$ and $u$ for its
degree-one element, so that $du=e_1-e_0$.  Also write
$e_a:\Disk{0}\to\Disk{1}$ for its endpoint functors and
$\epsilon:\Disk{1}\to\Disk{0}$ for the unique functor.  The projection
is
\[
 \pi_X=\id_X\otimes\epsilon:
 X\otimes\Disk{1}\longrightarrow X\otimes\Disk{0}\cong X.
\]
No interchange of the two tensor factors is used.

\subsection{Right cylinders and truncation}

Biclosedness supplies an $\omega$-category $H(Y)$ and natural
bijections
\begin{equation}\label{eq:right-internal-hom}
 \Hom_{\omCat}(X,H(Y))
   \cong\Hom_{\omCat}(X\otimes\Disk{1},Y).
\end{equation}
In the notation of \cite[\S A.18]{AraMaltsiniotis2020},
$H(Y)=\Hom_{\mathrm{oplax}}(\Disk{1},Y)$; the representing formula
\eqref{eq:right-internal-hom} fixes the convention.
Evaluation at $e_0,e_1$ and precomposition with $\pi_X$ give,
respectively, the endpoints and constant cylinders.
To extend cylinders through pushouts in $\nCat$, we will transpose
them to functors into $H(Y)$.  The next lemma shows that this
construction stays within $\nCat$ when the target $Y$ does.

\begin{lemma}[Internal hom and reflection]\label{lem:internal-truncation}
Let $n<\infty$ and let $Y$ be a strict $n$-category.  Then $H(Y)$ is a
strict $n$-category.  Consequently the reflection unit $X\to\trunc X$
induces a natural bijection
\begin{equation}\label{eq:truncated-cylinder-bijection}
 \Hom_{\omCat}((\trunc X)\otimes\Disk{1},Y)
 \ \xrightarrow{\ \cong\ }\
 \Hom_{\omCat}(X\otimes\Disk{1},Y).
\end{equation}
This bijection respects endpoints and restriction along functors in
the first variable.  For $n=\infty$ the same assertion holds with
$\trunc=\id$.
\end{lemma}

\begin{proof}
We prove the special case needed here directly; the general
internal-hom truncation assertion is also
\cite[Proposition~A.29]{AraMaltsiniotis2020}.
An $i$-cell of $H(Y)$ is a functor
$P:\Disk{i}\otimes\Disk{1}\to Y$.  Let
$s:\Disk{i-1}\to\Disk{i}$ be the source face inclusion, and let
$p:\Disk{i}\to\Disk{i-1}$ send the top $i$-cell to the identity of the
top $(i-1)$-cell.  Thus $ps=\id$ and $q=sp$ collapses
$\Disk{i}$ onto its source face.  Under
\eqref{eq:right-internal-hom}, the identity of the source of $P$ is
represented by $P(q\otimes\id)$.  We show that these functors agree
when $i>n$.

The atoms of $\Disk{i}\otimes\Disk{1}$ outside the source face
$\Disk{i-1}\otimes\Disk{1}$ come from exactly the following pairs:
the top $i$-generator or the target $(i-1)$-generator of $\Disk{i}$,
paired with $e_0,e_1$, or $u$.  Their dimensions are respectively
$i,i,i+1$ and $i-1,i-1,i$.  If $i>n+1$, all have dimension greater
than $n$.  If $i=n+1$, only the target $n$-generator paired with
$e_0$ or $e_1$ has dimension at most $n$.  At either endpoint, the
top $(n+1)$-cell maps under $P$ to an identity, so its source and
target $n$-cells have the same image.  These are precisely the two
equalities needed to compare $P$ with $P(q\otimes\id)$ on the
remaining degree-$n$ atoms.  All atoms in the source face are fixed
by $q\otimes\id$.

The two functors therefore agree on all atoms of dimension at most
$n$.  They agree on higher atoms as well: in an $n$-category the
image of a cell of dimension greater than $n$ must be the iterated
identity on its $n$-source, whose image has already been determined.
Generation by atoms proves $P=P(q\otimes\id)$.  Every cell of $H(Y)$
above $n$ is thus an identity.

Now transpose a cylinder to a functor $X\to H(Y)$.  Since $H(Y)$ is
an $n$-category, it factors uniquely through $\trunc X$.  Transposing
back proves \eqref{eq:truncated-cylinder-bijection}.  Naturality of
the reflection and of the adjunction proves the restriction
assertion.  The endpoint assertions follow either by naturality of
evaluation or by uniqueness of the factorization into $Y$.
\end{proof}

\begin{remark}\label{rem:omega-tensor-at-finite-n}
Even when $X$ and $Y$ are $n$-categories, the expression
$X\otimes\Disk{1}$ denotes their tensor in $\omCat$; it can have cells
above $n$.  A map from this tensor to $Y$ is nevertheless represented
by a map $X\to H(Y)$ \emph{in $\nCat$}.  This is the fact that permits
stationary cylinders to be glued by arbitrary pushouts in $\nCat$.
\end{remark}

\subsection{The Alexander--Whitney map}

For the Gray-cylinder-to-Street-nerve construction, see
Ara \cite[Sections~5.2.2--5.2.6, especially Proposition~5.2.3 and the
right-cylinder formula in~5.2.6]{Ara2023}.
The verification needed here is included below; the classical
Alexander--Whitney diagonal is credited separately in the proof.

\begin{lemma}[A right Gray cylinder gives a nerve homotopy]
\label{lem:gray-nerve}
There is a natural simplicial map
\[
 J_X:N_\infty X\times\Delta[1]\longrightarrow
           N_\infty(X\otimes\Disk{1})
\]
whose restriction at the endpoint $a$ is
$N_\infty(\id_X\otimes e_a)$, and for which
$N_\infty(\pi_X)J_X$ is the projection to $N_\infty X$.
Consequently a strict functor $P:X\otimes\Disk{1}\to Y$ induces a
simplicial homotopy between the nerves of its endpoint functors.
If a functor $a:A\to X$ satisfies
$P(a\otimes\id)=h\pi_A$ for some $h:A\to Y$, the homotopy restricts
along $N_\infty a$ to the constant homotopy of $N_\infty h$.
The same assertions apply to $\Street$ when $X,Y$ are $n$-categories.
\end{lemma}

\begin{proof}
For each $m$ define the classical Alexander--Whitney diagonal
\cite[IV.2, Remark~2.5]{GoerssJardine1999}
\begin{equation}\label{eq:aw}
 \operatorname{AW}_m[i_0,\ldots,i_p]
   =\sum_{a=0}^{p}
      [i_0,\ldots,i_a]\otimes[i_a,\ldots,i_p].
\end{equation}
It is a positive chain map $M[m]\to M[m]\otimes M[m]$.
For clarity, in the differential of the right side, deleting a vertex
strictly before or after the shared vertex gives the corresponding
term in $\operatorname{AW}_m d$, with coefficient $(-1)^j$ for
deletion of $i_j$.  The remaining terms cancel in pairs: deleting
the final vertex of the first factor at cut $a$ contributes
\[
 (-1)^a[i_0,\ldots,i_{a-1}]\otimes[i_a,\ldots,i_p],
\]
whereas deleting the initial vertex of the second factor at cut
$a-1$ contributes the same tensor with coefficient $(-1)^{a-1}$.
These cancellations use exactly the sign in
\eqref{eq:tensor-differential}.  The degree-zero assertion is
immediate.  Every summand in \eqref{eq:aw} is positive.
On vertices the augmentation is one; applying the augmentation to
the first or second tensor factor selects, respectively, only the
cut $a=0$ or $a=p$.  Thus both counits give the identity.

The diagonal is natural for every monotone ordinal map.  For a tuple
whose image is strictly increasing, this follows term by term.  If
the image has equal consecutive entries in positions $j,j+1$, then
at every cut $a\le j$ the second factor is degenerate, and at every
cut $a\ge j+1$ the first factor is degenerate.  Both sides of
naturality are therefore zero after normalization.  This verifies
degeneracies as well as faces.  Only this cosimplicial naturality is
asserted.

Apply $\nu$ and \eqref{eq:steiner-tensor} to obtain natural functors
$\delta_m:\Orient{m}\to\Orient{m}\otimes\Orient{m}$.
For an $m$-simplex $F:\Orient{m}\to X$ and a monotone map
$\theta:[m]\to[1]$, put
\[
 J_X(F,\theta)
   =(F\otimes\Orient{\theta})\delta_m,
\]
where $\Orient{\theta}:\Orient{m}\to\Orient{1}=\Disk{1}$ is the
cosimplicial structure map.  The naturality just checked makes
$J_X$ simplicial.  If $\theta$ is constant at $a$, this map factors
through the right counit and the endpoint $e_a$, giving the stated
endpoint formula.  Applying $\pi_X$ also gives the right counit,
so $N_\infty(\pi_X)J_X$ is the projection.  Postcomposing with
$N_\infty P$ proves the homotopy assertion, and naturality in $X$
together with the projection formula proves stationarity.
\end{proof}

\subsection{Finiteness and gluing in simplicial sets}

We use the Kan--Quillen model structure on $\sSet$.  Its cofibrations
are the monomorphisms, every object is cofibrant, and the model
structure is proper.  We also use the classical Kan replacement
$K\to\Ex^\infty K$, which is a weak equivalence and has Kan target.
These are simplicial homotopy-theory inputs; they do not presuppose
any model structure on strict higher categories
\cite[II.8.6 and III.4.8]{GoerssJardine1999}.

\begin{lemma}[Presentability and filtered colimits]\label{lem:filtered}
For every $1\le n\le\infty$, the category $\nCat$ is locally finitely
presentable, and filtered colimits are computed on cells.
Each $\trunc\Orient{m}$ is finitely presentable in $\nCat$.
Consequently $\Street$, and also $\Ex^r\Street$ for every finite
$r\ge0$, preserve filtered colimits.  Simplicial weak equivalences
are closed under filtered colimits of arrows, as are the maps
$f$ of $\nCat$ for which $\Street f$ is a weak equivalence.
No monomorphism hypothesis on transition maps is required.
\end{lemma}

\begin{proof}
Strict $n$-categories are models of a small finitary essentially
algebraic theory.  The sorts are the cell dimensions; the operations
are source, target, identity, and the binary compositions.  A
composition is defined by finitely many matching-boundary equations,
and every strict category axiom is an equation involving finitely
many cells.  In the $\omega$-case there are countably many sorts,
but all operations and equations still involve finite data.
Proposition~\ref{prop:strict-presentability} proves local finite
presentability, completeness, cocompleteness and a common regular
presentability bound directly from this signature. Its presentation
construction includes composites made possible by new boundary
identifications. The general essentially algebraic approach is developed
in \cite{AdamekRosicky1994,AdamekHebertRosicky1999}; for its finitary
partial-operation formulation, see also
\cite[pp.~52--53]{HebertAdamekRosicky2001}.

For the cellwise colimit assertion one can argue directly.  Form
the filtered colimit of the sets of cells in each dimension.
Any finite collection of representatives is represented at a common
stage.  If their boundaries agree in the colimit, those finitely many
equalities hold at some later common stage, where the required
composition is defined.  Define its colimit value by that
composition.  A further common stage proves independence of choices,
and each strict category law follows at a stage containing its
finitely many entries.  This constructs the colimit with its
universal property.  If all inputs are $n$-categories, every cell
above $n$ is already an identity, so the same is true of the colimit.

The oriental $\Orient{m}$ has a finite polygraphic presentation:
there is one generator for each nonempty strictly increasing tuple
in $[m]$, and its boundary is an expression in lower generators.
A map from $\Orient{m}$ into a filtered colimit is thus specified by
finitely many cell representatives satisfying finitely many
boundary equations.  These representatives and equations occur at
one stage.  Equality of two such maps is likewise detected at one
stage.  Hence $\Orient{m}$ is finitely presentable in $\omCat$.
The inclusion $\nCat\hookrightarrow\omCat$ preserves filtered
colimits by the preceding cellwise calculation, so reflection and
adjunction give finite presentability of $\trunc\Orient{m}$.
This proves the assertion for $\Street$ degreewise; compare also
\cite[Proposition~5.13(a)]{AraMaltsiniotis2014}.
Since
$(\Ex K)_m=\Hom_{\sSet}(\Sd\Delta[m],K)$ and
$\Sd\Delta[m]$ is a finite simplicial set, $\Ex$ preserves filtered
colimits too.

Here are details of the weak-equivalence assertion.  The functor
$Q=\Ex^\infty$ preserves filtered colimits, since $\Ex$ does and
colimits commute with colimits.  A filtered colimit of Kan complexes
is Kan: any horn has finitely many nondegenerate simplices, so it
and its compatibility equations occur at one stage and can be filled
there.  For a filtered diagram of Kan complexes $K_i$, components
commute with its colimit, since vertices and finite paths are finite
data.  If $x$ is a vertex of $\colim_iK_i$, choose a representative
$x_i\in K_i$.  The coslice of arrows $i\to j$ is filtered and its
projection to the indexing category is final; it gives a diagram
with compatible basepoints $x_j$.  For every $k\ge1$,
\[
 \colim_{(i\to j)}\pi_k(K_j,x_j)
       \ \cong\ \pi_k(\colim_jK_j,x).
\]
Indeed, use the finite pointed sphere
$\Delta[k]/\partial\Delta[k]$ and its finite reduced simplicial
cylinder to represent classes and pointed homotopies in a Kan
complex.  Their simplices and all required equalities, including
the basepoint equality, occur at a sufficiently late stage.  This
proves both surjectivity and injectivity of the displayed map,
including eventual equalities when transition maps are not
injective.  It follows that an objectwise weak equivalence between
filtered diagrams of Kan complexes induces a weak equivalence of
colimits.  Apply this to $Q$ of a diagram of arbitrary simplicial
weak equivalences, and use the natural weak equivalences $K\to QK$
and two-out-of-three.  Finally apply preservation of filtered
colimits by $\Street$.
\end{proof}

\begin{lemma}[Gluing with one monomorphic leg]
\label{lem:simplicial-gluing}
Consider a commutative map of spans of simplicial sets
\[
 \begin{tikzcd}[column sep=large]
 A \arrow[d,"a"'] &
 C \arrow[l] \arrow[r,hook,"i"] \arrow[d,"c"'] &
 B \arrow[d,"b"]\\
 A' &
 C' \arrow[l] \arrow[r,hook,"i'"'] &
 B'.
 \end{tikzcd}
\]
where $i$ and $i'$ are monomorphisms.
If $a,b,c$ are weak equivalences, then the canonical map
\[
 A\amalg_C B\longrightarrow A'\amalg_{C'}B'
\]
is a weak equivalence.  There is no condition on the maps $C\to A$
and $C'\to A'$.  Each of the displayed ordinary pushouts computes a
homotopy pushout.  In particular, cobase change of a weak equivalence
along a simplicial monomorphism is a weak equivalence.
\end{lemma}

\begin{proof}
This is the simplicial gluing lemma
\cite[II.8.8 and II.8.12]{GoerssJardine1999}; the following proof
keeps track of the specified cofibrations.  Put
\[
 E=B\amalg_C C'.
\]
By simplicial left properness, $B\to E$ is a weak equivalence,
since it is the pushout of $c$ along the monomorphism $i$.  The
induced map $E\to B'$ is consequently a weak equivalence by
two-out-of-three.  Both $C'\to E$ and $C'\to B'$ are cofibrations.
They are therefore cofibrant objects of the undercategory
$C'/\sSet$, with its inherited model structure.  Cobase change along
$C'\to A'$ is a left Quillen functor
\[
 C'/\sSet\longrightarrow A'/\sSet:
 \quad (C'\to T)\longmapsto(A'\to A'\amalg_{C'}T).
\]
Indeed, its right adjoint is restriction of the structure map,
which preserves fibrations and trivial fibrations.  It preserves
weak equivalences between cofibrant objects, so
\[
 A'\amalg_C B\ \cong\ A'\amalg_{C'}E
       \longrightarrow A'\amalg_{C'}B'
\]
is a weak equivalence.  Also $A\to A\amalg_CB$ is a monomorphism,
being a pushout of $i$.  Left properness applied to $a$ gives a
weak equivalence
$A\amalg_CB\to A'\amalg_CB$.  Composing the two proves the result.

To see the homotopy-pushout assertion, factor the possibly
nonmonomorphic map $C\to A$ as a cofibration $C\hookrightarrow
\widetilde A$ followed by a trivial fibration
$\widetilde A\to A$.  The span
$\widetilde A\leftarrow C\hookrightarrow B$ is a cofibrant
replacement of the original span for calculating the derived
pushout: its three objects are cofibrant and both legs are
cofibrations.  The already proved assertion identifies its pushout
with $A\amalg_CB$ up to weak equivalence.
For the final assertion, one may either use left properness directly
or apply the span assertion to
$X\leftarrow X\hookrightarrow Z$ and
$Y\leftarrow X\hookrightarrow Z$, with vertical maps
$X\to Y,\id_X,\id_Z$.
\end{proof}
